\documentclass[border=4pt]{standalone}
\usepackage{amsmath,amssymb,mathtools,xcolor,varwidth}
\definecolor{timeblue}{HTML}{1E6FE0}
\definecolor{meangreen}{HTML}{00A35A}
\definecolor{axisorange}{HTML}{E8770A}
\begin{document}
\begin{varwidth}{28em}
\large\bfseries
Here is the engine of the proof. By Prop. I the areal velocity is constant, so      \textcolor{timeblue}{the periodic time is the whole area $\pi a b$ divided by that steady rate}.      By Cor. Prop. XIV \textcolor{meangreen}{the lesser axis $b$ is the mean proportional between the greater axis      $a$ and the latus rectum $L$}, giving \textcolor{meangreen}{$b^2 = aL$}, whence $b \propto \sqrt{aL}$.      Since the force is as $1/SP^2$ (Prop. XI), the latus rectum enters only in its subduplicate ratio;      subducting that half-power leaves \textcolor{axisorange}{the sesquiplicate ratio $T \propto a^{3/2}$}, that is      \textcolor{axisorange}{$T^2 \propto a^3$}. Thus from ellipse, area, and inverse-square, Kepler's crown descends by pure geometry. \textit{Q.E.I.}
\end{varwidth}
\end{document}
